\documentclass[10pt,letterpaper]{article}
\usepackage[margin=.75in]{geometry}
\usepackage{amsmath,amssymb,graphicx,booktabs,xcolor,hyperref,fancyhdr}
\definecolor{navy}{HTML}{18384B}\definecolor{teal}{HTML}{168F9A}\definecolor{pale}{HTML}{EAF4F4}
\hypersetup{colorlinks=true,urlcolor=teal,linkcolor=teal,pdftitle={Two views of dispersive wave packets: a mock report}}
\pagestyle{fancy}\fancyhf{}\lhead{\small\color{navy}DISPERSIVE WAVE PACKETS}\rhead{\small\color{teal}Mock report / September 2026}\cfoot{\thepage}
\setlength{\headheight}{14pt}\setlength{\parindent}{0pt}\setlength{\parskip}{6pt}
\newcommand{\pic}[1]{\begin{center}\includegraphics[width=\linewidth]{#1}\end{center}}
\newcommand{\callout}[1]{\par\smallskip\colorbox{pale}{\parbox{.96\linewidth}{\vspace{5pt}#1\vspace{5pt}}}\par\smallskip}
\newcommand{\tagline}[1]{{\color{teal}\bfseries #1}\par}
\DeclareMathSizes{9}{10}{7}{7}\DeclareMathSizes{10}{10}{7}{7}\DeclareFontShape{OT1}{cmr}{bx}{n}{<->cmbx10}{}
\newcommand{\SelectedRows}{Random seed 33 & 0.003767 & 0.001491 & 60.4\% & 0.005556 \\
Two packets & 0.001433 & 0.000844 & 41.1\% & 0.003667 \\}

\begin{document}
\thispagestyle{empty}
\tagline{MOCK RESEARCH REPORT / SELECTED REVIEW FIGURES 5--8}
{\color{navy}\Huge Two views of dispersive\\[3pt] wave packets}\par
{\large A random envelope and a pair of overlapping packets in three spatial dimensions}\par
\vspace{6pt}
\textbf{Abstract.} We compare two artificial initial fields evolved by the same
homogeneous cubic dispersive operator: a seeded random packet and a pair of
phase-offset Gaussian packets. Three-dimensional level surfaces and planar slices
show how their shapes change between $t=0$ and $8$. Both sampled envelope peaks
fall, although their relative reductions differ. A supplementary time series
separates this visible spreading from the fractionally weighted diagnostic in the
proposed decay inequality. The results are finite-grid observations and support
an exploratory comparison; they do not identify an asymptotic exponent.
\pic{../output/figure_review/05_random_seed_33_volumes.png}
\textbf{Figure 5. Random seed 33: the envelope as a volume.} These are surfaces
of $|u|=0.3\max|u|$, using the full sampled maximum separately at each time.
All axes are spatial coordinates, and every panel uses the same $[-24,24]^3$
viewing cube. The magnitude is represented by the surface threshold, not by a
vertical elevation axis. The numerical maxima appear above the panels.

The initially compact central region becomes elongated and asymmetric. The
geometry changes while the sampled maximum drops from approximately $0.00377$
to $0.00149$. Because the threshold also falls, the growing visible region cannot
be interpreted as an increase of amplitude or as a quantitative measure of volume.
\callout{\textbf{Main observation.} The two selected fields illustrate different
routes to a lower peak: deformation of a random envelope and overlap of two
phase-offset packets. Their figures describe shape; the diagnostic comparison
on page 5 measures amplitude and its relation to the decay question.}
\vfill
{\small Dimensionless linear canonical model associated with the supplied Dysthe
research manuscript. Figure numbers 5--8 retain their identifiers from the review
collection. This draft is organized for discussion and figure selection.}

\newpage
\tagline{CASE A / LOOKING INSIDE A RANDOM ENVELOPE}
\section*{From a compact magnitude profile to an asymmetric packet}
The first input combines twelve complex plane waves under a Gaussian window:
\[
 f_{33}(X)=\frac{e^{-|X|^2/18}}{Z_{33}\sqrt{12}}
 \sum_{j=1}^{12}e^{i(q_j\cdot X+\phi_j)},\qquad X=(x,y,z).
\]
The wavevectors are sampled from independent centered normal components with
standard deviation $0.25$; phases are uniform on $[0,2\pi)$. NumPy seed 33 fixes
the realization. The positive constant $Z_{33}$ normalizes the sampled $L^1$ norm
to one. This is one reproducible realization, not a statistical ensemble average.
\pic{../output/figure_review/06_random_seed_33_slices_and_waves.png}
\textbf{Figure 6. Random seed 33: amplitude and oscillations.} Top panels show
$|u(x,0,z,t)|$. Their shared color scale divides by the largest displayed slice
magnitude across the three times. Bottom panels show the centerline envelope
$\pm|u(x,0,0,t)|$ and the illustrative carrier reconstruction
$\operatorname{Re}\{u(x,0,0,t)e^{3ix}\}$, divided by the same constant.

The fixed color scale makes the reduction of the bright central region visible.
The later slices also reveal weaker oscillatory structure away from the center.
Those lower-amplitude features are largely absent from the 30\% level surfaces
in Figure 5: a slice and a level surface answer different questions.

The dark oscillations remain inside the positive and negative envelope. Their
locations depend on the phase of the complex solution as well as on the added
carrier. The complex centerline is Fourier-interpolated from spacing $0.75$ to
$0.125$ for drawing; this makes the reconstruction smooth without adding simulated
Fourier modes. The carrier wavenumber 3 is an illustrative choice.
\callout{\textbf{Reading the magnitude.} It is dimensionless envelope amplitude,
not velocity. Neither these oscillations nor the 3D level surfaces represent a
calibrated water elevation. A physical surface reconstruction would require a
specified water-wave model and physical scales.}

\newpage
\tagline{CASE B / TWO PACKETS WITH A RELATIVE PHASE}
\section*{Overlap changes the combined shape}
The second input is deliberately simpler than the random realization:
\[
 f_{\rm pair}(X)=\frac{1}{Z_{\rm pair}}
 \left[e^{-((x+5)^2+y^2+z^2)/18}
       +0.7i\,e^{-((x-5)^2+y^2+z^2)/18}\right].
\]
The packet centers begin at $x=-5$ and $x=5$. Each has width 3; the second has
relative amplitude $0.7$ and a phase offset of $\pi/2$. As before, $Z_{\rm pair}$
is chosen using the grid quadrature so that the initial $L^1$ norm is one.
\pic{../output/figure_review/07_two_interfering_packets_volumes.png}
\textbf{Figure 7. Two packets: evolving 3D level surfaces.} The same 30\%-of-current-peak
rule and spatial viewing cube are used as in Figure 5. The separated initial
regions develop an overlapping, asymmetric shape. The plots show thresholded
magnitude geometry; they do not track individual fluid parcels.

The sampled peak decreases from $0.00143$ to $0.000844$, a smaller relative drop
than for the random packet over the same interval. Both inputs have $L^1=1$, but
they have different initial peaks, $L^2$ norms, and spectra. Equal $L^1$ normalization
therefore does not make their initial pointwise amplitudes equal.

\textbf{Why phase matters even in a linear model.} Let $u_1$ and $u_2$ denote the
evolved, normalized components, including the factor $0.7i$ in $u_2$. Linearity gives
\[
 u=u_1+u_2,\qquad
 |u|^2=|u_1|^2+|u_2|^2+2\operatorname{Re}(u_1\overline{u_2}).
\]
The last term is interference. Initially the Gaussian components have a quadrature
phase offset, so that term vanishes. Their subsequent complex evolution can make
it nonzero. The overlapping magnitude thus need not be the sum of the two separate
magnitudes.
\callout{\textbf{Interpretation.} This is linear superposition with interference.
The apparent joining of level surfaces is not evidence of nonlinear attraction,
soliton interaction, or irreversible merger.}

\newpage
\tagline{CASE B / CONNECTING THE VOLUME TO THE WAVE TRACE}
\section*{The slice explains what the level surface leaves out}
The planar and centerline views show the same two-packet computation from a more
familiar perspective. The two initial amplitude peaks become less distinct as
their evolved fields overlap, while the carrier oscillations remain visible.
\pic{../output/figure_review/08_two_interfering_packets_slices_and_waves.png}
\textbf{Figure 8. Two packets: fixed-scale slices and carrier reconstruction.}
The conventions are identical to Figure 6, but the normalization divisor is
computed separately for this case. Colors can be compared across times within
this figure; they cannot be used for an absolute amplitude comparison with
Figure 6. The printed full-volume peaks and the table on page 5 provide that comparison.

The upper panels retain the weaker structures below the threshold used in Figure 7.
The bottom row displays $\pm|u|$ as an outline rather than as an additional physical
wave. The dark curve combines the simulated phase with the illustrative carrier;
its crests are not numerical grid markers or tracked particles.

\textbf{What the comparison reveals.} The random realization begins with a more
concentrated peak and loses a larger fraction of that peak. The two-packet profile
spreads and overlaps while retaining a broader region of visible amplitude.
These differences persist despite identical evolution rules and initial $L^1$
normalization. Initial phase and spectral content are therefore essential parts
of the experiment definition.

\textbf{What the comparison cannot decide.} The figures do not determine a
universal decay constant, establish monotonicity for arbitrary data, or predict
which profile will have the larger compensated ratio at much later times. A
centerline may also miss an off-axis maximum. The following diagnostics therefore
search the entire simulated volume rather than one selected slice.
\callout{\textbf{A useful pairing for the final report.} Figures 5 and 7 establish
the three-dimensional geometry. Figures 6 and 8 explain the amplitude and phase
behind that geometry. Keeping each volume next to its companion slice avoids
asking one visualization to communicate both shape and decay.}

\newpage
\tagline{SUPPLEMENT / CONNECTING THE FIGURES TO THE DECAY QUESTION}
\section*{A falling envelope peak is only one diagnostic}
For the same two inputs and the same grid, we evaluate 17 times from 0 to 8.
The weighted field and compensated ratio are
\[
 W(t)=|D_x|^{1/2}S(t)f,\qquad
 R(t;f)=\frac{t^{7/6}\|W(t)\|_\infty}{\|f\|_1}.
\]
These are different from the unweighted envelope magnitudes shown in Figures 5--8.
The ratio is displayed only for $t>0$.
\pic{../results/selected_packets/comparison.pdf}
\textbf{Figure S1. Same inputs, four complementary measurements.} Maxima are
sampled over the full periodic volume. The upper right panel removes the difference
in initial peak height. The lower right panel measures the quantity associated
with the manuscript's $\alpha=1/2$, $p=\infty$ bound.
\begin{center}\small\begin{tabular}{lrrrr}\toprule
Input & Peak at 0 & Peak at 8 & Peak reduction & $R(8)$\\\midrule
\SelectedRows\bottomrule\end{tabular}\end{center}
\callout{\textbf{The relevant distinction.} A decreasing unweighted peak can
coexist with an increasing compensated ratio. The proposed inequality asks
whether $R(t;f)$ stays below an input-independent constant, not whether every
finite curve decreases.}
All curves on this page are measurements on one finite grid. Their purpose is to
connect the selected visuals to the research question; they are not a new
convergence study or an asymptotic fit.

\newpage
\tagline{METHODS, RELIABILITY AND CONTEXT}
\section*{What was computed, and how to use this draft}
\textbf{Evolution.} We work directly in the homogeneous canonical coordinates of
the supplied manuscript, with
\[
 \widehat u(k,t)=e^{itw(k)}\widehat f(k),\qquad
 w(k)=k_x^3-3k_x(k_y^2+k_z^2).
\]
Equivalently, $u_t+u_{xxx}-3u_{xyy}-3u_{xzz}=0$. The Fourier convention is
$\widehat f(k)=\int f(X)e^{-ik\cdot X}\,dX$. This is the linear propagator used
in the homogeneous weighted estimate; it is not the full nonlinear Dysthe equation.

\begin{center}\begin{tabular}{ll}\toprule
Setting & Value\\\midrule
Computational domain & Periodic cube $[-72,72)^3$\\
Grid & $192^3$, spacing $0.75$\\
Displayed region & $[-24,24]^3$; slices at $y=0$\\
Snapshot times & $0,4,8$\\
Supplementary time series & $0,0.5,\ldots,8$\\
Arithmetic & Torch CPU, complex128\\
Initial normalization & Discrete $L^1=1$ for each input\\
Time evolution & Exact Fourier phase for retained modes\\\bottomrule
\end{tabular}\end{center}
\textbf{Checks and their limits.} The new diagnostic series reproduces the three
selected snapshot times. Integrated squared magnitude is conserved to numerical
precision: the relative $L^2$ changes are below $10^{-11}$. Conservation checks
consistency but cannot establish an accurate whole-space maximum or detect every
phase error. No separate domain enlargement or mesh refinement was performed for
this selected-packet series.

At $t=8$, the maximum sampled amplitude on the displayed cube's boundary is
approximately 16.4\% of the full peak for seed 33 and 6.0\% for the two-packet case.
Thus the 30\% surfaces do not touch that sampled boundary, but weaker tails are
cropped from view. This is a display check, not a test of the computational boundary.
The strong $t=256$ audit in the larger project applies to other Gaussian-based
inputs and must not be transferred to these two cases.

\textbf{Research context.} The supplied manuscript by Mosincat, Pilod and Saut,
\emph{On Third Order Nonlinear Schr\"odinger Equations} (January 22, 2026),
Section 6, provides the canonical symbols and weighted estimates. Its proof-phase
convention is used here. Ben-Artzi, Koch and Saut's
\href{https://doi.org/10.1081/PDE-120025491}{two-dimensional cubic dispersion study}
(2003) motivates attention to spectral geometry. For a physical free-surface model,
Guyenne, Kairzhan and Sulem's
\href{https://arxiv.org/abs/2108.10822}{Hamiltonian Dysthe study} (2022) offers
surface reconstruction and Euler comparisons; its fluid domain has two horizontal
surface coordinates, unlike the three independent coordinates here.

\textbf{Suggested scientific follow-up.} First compare domain and grid changes
for these exact two recipes and the weighted peaks. Then extend their observation
window while tracking the global candidate maximum. Repeating seed 33 with more
seeds would support an ensemble comparison; one realization cannot characterize
random initial data in general.

\begingroup\raggedright\textbf{Reproduction.} Run \path{make_figure_review.py} for the selected figures.
Run \path{report/selected_packet_diagnostics.py} for Figure S1, its table, and the
JSON records in \path{results/selected_packets/}. Compile
\path{report/selected_packets_mock.tex} to rebuild this draft.\par\endgroup
\end{document}
